\section{Hai tiên đề cho attribution}
\label{sec:ch05-hai-tien-de}

\index{độ nhạy}%
Hai khiếm khuyết của mục trước, bão hòa và đơn vị đo, cho thấy cần một tiêu
chuẩn nằm ngoài trực giác.
Sundararajan, Taly và Yan bắt đầu bằng hai yêu cầu đối với một lời giải thích
gán đóng góp cho đầu vào~\autocite{p11ig}. Yêu cầu thứ nhất là \tn{độ nhạy}{sensitivity}:
nếu đầu vào và baseline chỉ khác ở một đặc trưng mà đầu ra khác nhau, đặc trưng
đó phải nhận đóng góp khác không. Điều kiện này loại trường hợp hàm đã đổi trên
đường đi nhưng phẳng tại đầu vào.

\index{bất biến theo cài đặt}%
Yêu cầu thứ hai là \tn{bất biến theo cài đặt}{implementation invariance}. Hai
mạng có kiến trúc hoặc tham số trung gian khác nhau nhưng tính đúng cùng một
hàm phải cho cùng lời giải thích. Attribution nhằm nói về quan hệ giữa đầu vào
và đầu ra, nên nó không nên thay đổi chỉ vì người triển khai viết lại cùng một
phép tính. Gradient thỏa yêu cầu này nhờ quy tắc dây chuyền, vì nó chỉ phụ thuộc vào hàm
mà mạng biểu diễn. Những phương pháp sửa lại quy tắc lan truyền ngược, gán cho
mỗi đơn vị của mạng một điểm số rồi truyền ngược điểm số ấy theo một quy tắc
riêng thay vì theo đạo hàm, thì không: bài báo nêu DeepLift và LRP làm ví dụ,
vì điểm số của chúng phụ thuộc vào cách mạng được viết chứ không chỉ vào hàm
mà mạng tính~\autocite{p11ig}.

Hai tiên đề không chứng minh một bản đồ là trung thực theo nghĩa nhân quả. Chúng
chỉ loại những lời giải thích mâu thuẫn với hai trực giác tối thiểu về đóng góp.
Chúng cũng không chọn baseline. Chính vì thế, phần còn lại của chương tách hai
việc thường bị gộp lại: phép tính attribution có thỏa các điều kiện hình thức
hay không, và đối chiếu với baseline có mang nghĩa đúng với bài toán hay không.

